% Owner: boundary_d.  Separate from the E-owned master source.
% These auxiliary assertions concern one-sided smooth specializations.
% The full deformation-chain conclusion is proved in E.2.22--E.2.24.
\subsubsection*{Smooth specializations of the canonical pencil}

\begin{proposition}[E.2.5 (Proved): a fractional canonical pencil at a smooth limit]
Let $\pi:\mathcal X\to\Delta$ be a proper holomorphic submersion
whose fibres are diffeomorphic to $S^6$.  Suppose that for every
$t\ne0$ the fibre is one of the members of E.2.1, with the same
coprime multiplicities $(3k,4l)$.  Set
\[
 A=12kl,\qquad D=kl-3k-4l\ne0.
\]
There is a unique holomorphic line bundle $\mathcal L$ on
$\mathcal X$, up to isomorphism, satisfying
\[
 \mathcal L^{D}=K_{\mathcal X/\Delta}^{A}.
\]
Its restriction to $X_t$ for $t\ne0$ is the coarse pencil
$f_t^*\mathcal O_{\mathbb P^1}(1)$.  The coherent sheaf
$\pi_*\mathcal L$ is locally free of rank two, and its central
fibre injects into $H^0(X_0,\mathcal L_0)$.

The central manifold has algebraic dimension one and a
holomorphic algebraic reduction $g_0:X_0\to\mathbb P^1$ with
connected fibres.  If $s_0,s_1$ are the two central sections
obtained from a local frame of $\pi_*\mathcal L$, there are a
nonconstant rational map $R:\mathbb P^1\to\mathbb P^1$ of
degree $e\ge1$ and an effective Cartier divisor $E_0$ such that
\[
 [s_0:s_1]=R\circ g_0,\qquad
 \mathcal L_0\simeq g_0^*\mathcal O(e)\otimes\mathcal O(E_0).
\]
Here $[s_0:s_1]$ denotes the holomorphic extension of the ratio;
the displayed divisor is its common fixed divisor.  The
statement does not assert that $E_0=0$.
\end{proposition}
\begin{proof}
By Ehresmann's theorem the smooth bundle over the contractible
disc is differentiably trivial.  Consequently
$H^1(\mathcal X,\mathbb Z)=H^2(\mathcal X,\mathbb Z)=0$.
The exponential sequence identifies
\[
 \operatorname{Pic}(\mathcal X)=H^1(\mathcal X,\mathcal O_{\mathcal X})
\]
as groups.  This is a complex vector space, so multiplication by
the nonzero integer $D$ is an automorphism.  This proves the
existence and uniqueness of the \emph{relative} root, including
when $D<0$.  The same argument on each fibre makes its Picard
group uniquely divisible.  E.2.1 gives
$K_{X_t}^{A}=(f_t^*\mathcal O(1))^{D}$, so uniqueness identifies
$\mathcal L_t$ with this actual pencil, not just with a numerical
class.

Proper direct-image coherence gives $\mathcal V=\pi_*\mathcal L$.
It is torsion-free: a local section killed by a nonzero function
of $t$ vanishes on the inverse image of an open set, hence
vanishes as a section of the invertible sheaf on the smooth
total space.  A coherent torsion-free sheaf on a smooth curve
is locally free.  Its rank is two because the noncentral
algebraic reductions have connected fibres.  Applying $\pi_*$
to
\[
 0\longrightarrow\mathcal L\xrightarrow{\,t\,}\mathcal L
 \longrightarrow i_*\mathcal L_0\longrightarrow0
\]
shows directly that $\mathcal V/t\mathcal V$ injects into
$H^0(X_0,\mathcal L_0)$.  Thus the two central sections are
linearly independent even if central cohomology jumps; no
surjectivity of base change has been assumed.

We use precisely the non-holomorphic-meromorphic-map theorem
of \citet[Theorem~2.1]{CDP2020D}, also written out in
\citet[Theorem~3.1 and Corollary~3.2]{LRS2018D}.
For a compact complex threefold with $b_2=0$, a nonconstant
meromorphic map to a curve which fails to be holomorphic
would force $c_3=0$.  Since $c_3(X_0)[X_0]=2$, every
meromorphic function on $X_0$ extends as a holomorphic map
to $\mathbb P^1$.  This argument uses neither the different
holomorphic-fibration assertion of CDP, Theorem~2.2, nor its
Corollary~2.3.  Those assertions cannot be applied here.

In particular the ratio of $s_0,s_1$ gives a nonconstant
holomorphic map $h_0:X_0\to\mathbb P^1$.  Also $a(X_0)\le1$:
otherwise two algebraically independent meromorphic functions
would give a holomorphic map $X_0\to\mathbb P^1\times\mathbb P^1$
with two-dimensional image.  At a point of differential rank
two, after choosing finite affine target coordinates $u,v$
centred at its image, the global meromorphic function $u/v$
would have the local indeterminacy of the quotient of two
independent coordinates.  It cannot extend holomorphically
there.  This contradicts the preceding theorem.  Hence
$a(X_0)=1$.

Take the Stein factorization $h_0=R\circ g_0$.  Its intermediate
space is a smooth compact curve.  Connectedness of the fibres
injects its $H^1(\ ,\mathbb Z)$ into $H^1(X_0,\mathbb Z)=0$
by the degree-one Leray sequence, so it is $\mathbb P^1$.
Every meromorphic function on $X_0$ factors rationally through
$g_0$: the image of its holomorphic map together with $g_0$
is an analytic curve in $\mathbb P^1\times\mathbb P^1$;
normalizing this curve gives a finite factor of $g_0$, whose
degree is one because $g_0$ has connected fibres.  Thus $g_0$
is the algebraic reduction, and $R$ has some degree $e\ge1$.

Finally, on a target chart where the first homogeneous
coordinate of $h_0$ is nonzero, $s_1=h_0s_0$ in a local
trivialization, and the common ideal is generated by $s_0$.
The analogous assertion holds on the other chart.  These local
generators define a nonzero section of
$\mathcal L_0\otimes h_0^*\mathcal O(-1)$ and an effective
Cartier zero divisor $E_0$.  This gives the stated factorization
of $\mathcal L_0$ with no residual codimension-two base locus.
\end{proof}

\begin{proposition}[E.2.6 (Proved): positive canonical slope preserves the fibration]
Under the hypotheses of E.2.5, suppose $D>0$.  Then $E_0=0$
and $e=1$.  After shrinking the disc and choosing a frame of
$\pi_*\mathcal L$, its evaluation defines a proper flat
holomorphic map
\[
 g:\mathcal X\longrightarrow\mathbb P^1\times\Delta
\]
which is the intrinsic algebraic reduction on every fibre.
Its smooth general central fibre is a complex two-torus.

In particular, two families of the type in E.2.5 with positive
canonical slope cannot have the same smooth central manifold
unless their multiplicity pairs agree.  A positive-slope and
a negative-slope family cannot have a common smooth central
manifold either.  This proposition concerns a common smooth
specialization.  The additional propagation and closure
arguments for arbitrary connected chains are E.2.22--E.2.24.
\end{proposition}
\begin{proof}
Choose a very general smooth fibre $F$ of $g_0$, avoiding the
images of all vertical components of $E_0$ and the jumping
sets for the pluricanonical direct images below.  Its normal
bundle in $X_0$ is trivial.  Adjunction gives
$K_F=K_{X_0}|_F$, and $c_1(K_F)=0$ integrally because
$H^2(X_0,\mathbb Z)=0$.  In particular $F$ is minimal: a
$(-1)$-curve would have canonical intersection $-1$.
Restriction of the root identity in E.2.5 gives
\[
 K_F^{A}\simeq\mathcal O_F(D E_0|_F).
\]
The right hand side is effective since $D>0$, so
$\kappa(F)\ge0$.

We next show $\kappa(F)\le0$, without assuming that the fibres
of an algebraic reduction have algebraic dimension zero.
If $\kappa(F)>0$, some $r>0$ has $h^0(F,K_F^r)\ge2$.
Our very-general choice ensures that this is the generic rank
of the coherent sheaf $g_{0*}K_{X_0}^r$.  A sufficiently
positive twist of this sheaf on $\mathbb P^1$ is globally
generated.  Two of its global sections can therefore be
chosen to be linearly independent on a general fibre.  By
projection formula they are sections of
$K_{X_0}^r\otimes g_0^*\mathcal O(N)$; their ratio is a
meromorphic function on $X_0$ which is nonconstant on $F$.
This contradicts the algebraic-reduction assertion of E.2.5.
The choice of a very general fibre is legitimate: for each
$r$ the exceptional set for the generic cohomology rank is
a proper analytic subset of the base curve, and there are
only countably many $r$.

Thus $\kappa(F)=0$.  The classification of minimal compact
complex surfaces makes $K_F$ torsion; the relevant consequence
is also explicitly recorded in the proof of
\citet[Proposition~3.1.2]{AbashevaDeev2024AB}.
An effective divisor whose line bundle is torsion is zero:
raising its canonical section to a torsion-killing power
gives a holomorphic function on a compact connected manifold,
and that function is a nonzero constant.  Therefore
$E_0|_F=0$.  Every horizontal component of $E_0$ would meet
such a general fibre in a nonempty effective divisor.
Consequently $E_0$ has only vertical components, supported
over finitely many points of $\mathbb P^1$.

We now use connectedness of actual divisor limits to remove
both this vertical fixed part and the possible degree $e$.
Choose a linear combination $s=as_0+bs_1$ for which its
moving central divisor is the union of the $e$ distinct
smooth fibres over $R^{-1}([-b:a])$, all disjoint from the
support of $E_0$.  Also choose it to be a regular fibre on
one specified noncentral member.  Both requirements omit
only finitely many choices in the projective line of
combinations.  Extend it by the corresponding combination
of the frame of $\pi_*\mathcal L$.

Its total zero divisor $\mathcal D\subset\mathcal X$ is a
relative effective Cartier divisor: the section is not
identically zero on any fibre, by independence of the two
sections, including at zero.  In local coordinates, $t$ is
a non-zero-divisor modulo its defining equation.  Hence
$\mathcal D\to\Delta$ is proper and flat.  Its general
noncentral fibre is a smooth connected torus.  There are no
vertical irreducible components of $\mathcal D$, since such
a component would have to be the whole irreducible
threefold $X_t$.  It follows that $\mathcal D$ is reduced
and irreducible: its unique generic component has multiplicity
one, and a Cartier divisor in a smooth manifold has no
embedded components.  Normalize $\mathcal D$ and take the
Stein factorization over $\Delta$.  The finite factor has
degree one over the punctured generic locus; normality of
the disc makes it the disc itself.  Its central fibre, and
therefore the image $\mathcal D_0$, is connected.

But the support of
\[
 \mathcal D_0=E_0+g_0^*R^*([-b:a])
\]
is a disjoint union of the $e$ chosen smooth connected fibres
and, if nonempty, the support of $E_0$ over other base points.
Connectedness forces $E_0=0$ and $e=1$.

The central two sections now generate $\mathcal L_0$ at
every point.  The noncentral two sections already generate
their pencils, so the evaluation
$\pi^*\pi_*\mathcal L\to\mathcal L$ is surjective after
shrinking.  It defines the asserted map $g$.  It is proper
because $\pi$ is proper.  On each $X_t$ the nonconstant
map to the smooth curve $\mathbb P^1$ is flat: its local
structure sheaf is torsion-free over the local discrete
valuation ring of the curve.  The total source is flat
over $t$, so the local criterion for flatness applied to
the parameter $t-t_0$ gives flatness over $(z,t)$.
Equivalently, the local fibre equations form a regular
sequence in the smooth source.  At a regular
point of the central base curve, properness makes $g$ a
proper submersion on a neighbourhood of the corresponding
compact fibre.  Ehresmann's theorem identifies that central
surface differentiably with nearby complex two-tori.
The classification of compact complex surfaces, or its
special case for the differentiable four-torus, then makes
the central surface a complex two-torus as well.

Suppose now two positive-slope families have the same
central manifold $X_0$.  Its algebraic reduction is intrinsic
up to $\operatorname{PGL}_2$, so the preceding result makes
both central pencil lines the same non-torsion line
$L=g_0^*\mathcal O(1)$.  With indices $i=1,2$ the identities
$K_{X_0}^{A_i}=L^{D_i}$ give
$L^{D_1A_2-D_2A_1}=\mathcal O$.  The line $L$ is non-torsion
because a positive power has a nonzero section with zeros.
Thus $D_1/A_1=D_2/A_2$, and the exact recovery of the ordered
pair from this slope in E.2.3 gives the assertion.

Finally, a common smooth limit from opposite signs would
have a nonzero section of a positive power of $K_{X_0}$
and of a negative power, by upper semicontinuity (or by
the same injective direct-image argument).  Raising them
to a common power, their nonzero product is a constant;
both are nowhere zero.  Hence $K_{X_0}$ is torsion.
The exponential sequence and $H^1(X_0,\mathbb Z)=
H^2(X_0,\mathbb Z)=0$ make $\operatorname{Pic}(X_0)$
torsion-free, so $K_{X_0}$ would be trivial.  Its nonzero
holomorphic three-form $\Omega$ is closed and exact because
$H^3(X_0,\mathbb C)=0$.  Stokes' theorem would give
$\int_{X_0}\Omega\wedge\overline\Omega=0$, whereas its
appropriate constant multiple is a positive volume form.
This contradiction excludes the common limit.
\end{proof}

\begin{remark}[Scope after E.2.6]
E.2.5 proves that the central meromorphic reduction is
holomorphic, and E.2.6 removes its fixed divisor and degree
loss for $D>0$.  The positive-slope proof alone does not
exclude the effective antipluricanonical divisor allowed
by the restriction identity when $D<0$.  E.2.8--E.2.10
remove that possibility by the horizontal-divisor vanishing
argument.  E.2.13 and E.2.16--E.2.21 then recover the finite
charts, cusp, and actual smooth completion.  The separate
generic-disc propagation and closure arguments in
E.2.22--E.2.24 prove E.2.4 for arbitrary analytic bases,
deformation chains, and continuous complex structures.
The present proposition retains its stated one-sided,
positive-slope scope.
\end{remark}
